% Section 3, failure B: G' remains an accepted prefix of old G; D forks.
\begin{tikzpicture}
\path[use as bounding box] (-.1,-.1) rectangle (23.0,16.8);
\node[heading] at (11.45,16.2) {故障二：G 继续增长，D 在窗口到期后与 G 分叉};
\node[note,anchor=west] at (.5,15.4) {共同起点：$R=115$，$G=108$，$G'=D=105$。新的批次交付失败；此前已接纳的数据仍可取回。};
\node[subheading,anchor=west,text=orange] at (.5,14.55) {阶段 1：G → $G'$ 受阻；R 与 G 继续增长，Inbox cursor 停住};
\node[box,fill=regionA,text width=3.8cm,minimum height=1.5cm] (r) at (2.7,13.15) {\textbf{R：区域结果}\\继续生产，覆盖 G};
\node[box,fill=regionA,text width=3.8cm,minimum height=1.5cm] (g) at (8.5,13.15) {\textbf{G：全局链}\\108 $\to$ 200（示例）};
\node[box,fill=regularB,draw=orange,text width=3.8cm,minimum height=1.5cm] (gp) at (14.3,13.15) {\textbf{$G'$：Inbox cursor}\\停在 105};
\node[box,fill=native,draw=violet,text width=3.8cm,minimum height=1.5cm] (d) at (20.1,13.15) {\textbf{D：派生链}\\到 105，等待窗口};
\draw[flow] (r) -- (g);
\draw[ref,draw=rose] (g.east) -- (gp.west);
\draw[rose,line width=1.5pt] (11.23,12.98) -- (11.57,13.32) (11.23,13.32) -- (11.57,12.98);
\node[tag,text=rose] at (11.4,13.75) {受阻};
\draw[flow] (gp) -- (d);
\node[note,anchor=west] at (.5,11.95) {Batcher、Certifier 或新的 DA 交付失败。开始时四条链仍兼容，只是 G 与 $G'$ 的进度差越来越大。};
\node[subheading,anchor=west,text=violet] at (.5,11.15) {阶段 2：106 的派生窗口到期；D 采用 deposit-only，旧 G 仍保留普通块};
\foreach \yy/\label in {9.8/R,8.05/G,6.3/{G'},4.55/D}{\node[rowlabel] at (2.15,\yy) {$\label$};}
\node[box,kept,fill=settled,text width=4.7cm] (br) at (5.2,9.8) {\textbf{相同前缀：至 105}};
\node[box,kept,fill=settled,text width=4.7cm] (bg) at (5.2,8.05) {\textbf{相同前缀：至 105}};
\node[box,draw=orange,fill=regularB,text width=4.7cm] (bp) at (5.2,6.3) {\textbf{cursor = 105}\\只记录接纳进度};
\node[box,kept,fill=settled,text width=4.7cm] (bd) at (5.2,4.55) {\textbf{相同前缀：至 105}};
\node[box,dead,text width=5.0cm] (oldr) at (12.2,9.8) {\textbf{旧生产后缀：106--200\ldots}\\普通块，无法接回 D};
\node[box,draw=blue,fill=regionA,text width=5.0cm] (oldg) at (12.2,8.05) {\textbf{旧 G：106--200 普通块}\\有全局确认，未按时交付};
\draw[flow] (br) -- (oldr); \draw[flow] (bg) -- (oldg);
\node[note,anchor=west,text width=5.1cm] at (9.4,6.3) {$G'$ 仍只接纳到 105；\\D 的补位不推进 cursor。};
\node[box,draw=violet,fill=native,text width=5.0cm] (dep) at (12.2,4.55) {\textbf{106--150：deposit-only}\\由派生窗口回退确定};
\node[box,draw=violet,fill=native,text width=4.05cm] (later) at (19.4,4.55) {\textbf{151\ldots}\\D 继续按窗口补位};
\draw[flow,draw=violet] (bd) -- (dep); \draw[flow,draw=violet] (dep) -- (later);
\draw[rose,line width=.9pt] (15.1,8.05) -- (15.55,8.05) -- (15.55,4.55) -- (15.1,4.55);
\node[note,text=rose,text width=5.4cm,anchor=west] at (16.1,7.3) {\textbf{第一个不同的块：106}\\G 是普通块；\\D 已确定为 deposit-only。};
\node[card,fill=rose!3,draw=rose,text width=20.5cm] at (11.45,2.55) {
  \textbf{结果：$G'$ 仍是旧 G 的接纳前缀，但 D 与 G、旧 R 都已分叉。}\\
  G 的高度领先，不代表 D 是它的前缀。已经结算的共同前缀只到 105，旧 G 的后缀要在 RESUME 时对齐。
};
\node[note,anchor=west,text width=21.4cm] at (.5,1.15) {四行是高度顺序的示意；D 的窗口补位逐步发生。G 不会在分叉被发现时自动撤销，旧生产由后文的 Kill 停止。};
\node[note,anchor=west,text width=21.4cm] at (.5,.45) {本例故障发生在新批次交付；若此前已接纳的 DA 数据取不回，D 会等待，不能靠 deposit-only 跳过。};
\end{tikzpicture}
